\documentclass{article}

\usepackage{array}
\usepackage{etoolbox}
\usepackage{fancyhdr}
\usepackage{geometry} 
\usepackage{graphicx}
\usepackage{soul}
\usepackage{titling}

%%%%%%%%%%%%%%%%%%%%%%%%%%%%%%%%%%%%%%%%%%%%%%%%%%%%%%%%%%%%
% BEGIN METADATA: Edit the following as appropriate
%%%%%%%%%%%%%%%%%%%%%%%%%%%%%%%%%%%%%%%%%%%%%%%%%%%%%%%%%%%%

\title{Project Title}  % the title of your project
\newcommand\shorttitle{\thetitle}  % if needed: a shorter title for the document header
% Team members.
\newcommand\firstname{Ifrah Chishti}  % full name
\newcommand\firstid{ic08351}         % ID, e.g. xy01234
\newcommand\secondname{Naaseh Sajid} % full name
\newcommand\secondid{ms08085}        % ID, e.g. xy01234
\newcommand\thirdname{Aina Shakeel}  % full name
\newcommand\thirdid{as08430}         % ID, e.g. xy01234
% Uncomment the rows for the next 2 students if and as needed.
\newcommand\fourthname{Zain Hatim} % full name
\newcommand\fourthid{zh08343}        % ID, e.g. xy01234
% \newcommand\fifthname{Student 5}  % full name
% \newcommand\fifthid{id05}         % ID, e.g. xy01234

%%%%%%%%%%%%%%%%%%%%%%%%%%%%%%%%%%%%%%%%%%%%%%%%%%%%%%%%%%%%
% END METADATA: Do not edit the preamble any further.
%%%%%%%%%%%%%%%%%%%%%%%%%%%%%%%%%%%%%%%%%%%%%%%%%%%%%%%%%%%%

\pagestyle{fancy}
\lhead{Kaavish Proposal}
\chead{\shorttitle}
\rhead{Fall 2025}
\cfoot{Page \thepage}
\renewcommand{\footrulewidth}{0.4pt}

\newenvironment{instruction}{\itshape}{}

\begin{document}

% Cover page.
\begin{titlepage}

\center % Center everything on the page
 
%----------------------------------------------------------------------------------------
%	HEADING SECTIONS
%----------------------------------------------------------------------------------------

\textsc{
  {\LARGE \bf \thetitle}\\\bigskip\bigskip % Your Project Title
  {\large
    Kaavish Project Proposal\\\bigskip
    By}
}\\\bigskip 

%----------------------------------------------------------------------------------------
%	AUTHOR SECTION
%----------------------------------------------------------------------------------------

{\large
  \begin{tabular}{ll}
    \firstname & (\firstid@st.habib.edu.pk) \\
    \secondname & (\secondid@st.habib.edu.pk) \\
    \thirdname & (\thirdid@st.habib.edu.pk) \\
    \ifdef{\fourthname}{\fourthname & (\fourthid@st.habib.edu.pk) \\}{}
    \ifdef{\fifthname}{\fifthname & (\fifthid@st.habib.edu.pk) \\}{}
  \end{tabular}
}
\bigskip\bigskip\bigskip

{\large \today}\\\bigskip\bigskip

\includegraphics[height=5cm]{HU_logo}\\\bigskip
 
%----------------------------------------------------------------------------------------
{\large
  In partial fulfillment of the requirement for \\\medskip
Bachelor of Science \\\medskip
Computer Science
}\\\bigskip\bigskip\bigskip

{\large
  \textsc{
    Dhanani School of Science and Engineering\\\bigskip
    Habib University\\\bigskip 
    Fall 2023
  }\\\bigskip\bigskip 
  Copyright @ 2023 Habib University
}

\end{titlepage}


%%%%%%%%%%%%%%%%%%%%%%%%%%%%%%%%%%%%%%%%%%%%%%%%%%%%%%%%%%%%
% DATA: Populate the rest of the document as instructed.
%%%%%%%%%%%%%%%%%%%%%%%%%%%%%%%%%%%%%%%%%%%%%%%%%%%%%%%%%%%%
\section{Abstract}
\begin{instruction}
    CodeSchool is a global coding literacy platform for children aged 6 - 16, delivering more than 2,000 classes across 20 countries. With a mission to make coding a universal literacy, the organization combines an engaging, game-based curriculum with trained instructors who can inspire the next generation of digital creators. Much like how Uber seamlessly connects riders with drivers, CodeSchool envisions a platform that can intelligently connect students with the right instructor, at the right time, using the right curriculum—whether classes are delivered online or onsite. To achieve this vision at scale, the platform must evolve from a system dependent on manual operations into one that is automated, intelligent, and easy to use.

    Currently, managing the lifecycle of classes at CodeSchool is a complex, time-consuming process. Instructors and administrators rely on spreadsheets to schedule classes, assign students, generate Zoom links, and track progress. Evaluations of teaching quality are also carried out manually, with the team watching recorded sessions and giving scores based on observations. While effective at a smaller scale, this system becomes increasingly difficult to manage as CodeSchool grows, leaving room for scheduling conflicts, inefficiencies, and limited insight into teaching quality. These challenges risk slowing down the organization's ability to expand while maintaining its high standard of student learning outcomes.

    The proposed solution is a web-based platform designed to act as a central instructor portal—streamlining operations, automating repetitive tasks, and enabling data-driven insights. At its core, the platform addresses three key areas: class scheduling, student-instructor matching, and performance evaluation.

    First, the platform will automate class scheduling. By factoring in instructor availability, coding language expertise, and student age and proficiency, it will generate conflict-free schedules. This not only saves administrative effort but also ensures that classes run smoothly without last-minute rescheduling.

    Second, the platform will transform how students are assigned to instructors. Instead of a manual process that relies on human judgment, the system will use a recommendation model inspired by machine learning. Much like matching services in other industries, this model will prioritize placing students with instructors best suited to their learning needs—whether that is based on teaching style, technical strengths, or experience with specific age groups. This will also follow a round robin model to ensure fair distribution of students among instructors, allowing instructors to experience teaching different age groups and fostering a more versatile teaching approach. The result is a more personalized and effective learning experience for every student.

    % Third, to support the logistics of virtual teaching, the platform will integrate with the Zoom API to automatically generate and distribute meeting links for each scheduled class. This eliminates another manual step and ensures that both students and instructors always have seamless access to their sessions.

    Finally, the platform will introduce a more advanced approach to instructor performance analysis. Rather than relying solely on manual reviews, it will leverage automated transcription and analysis of Zoom recordings. Using metrics such as clarity, pacing, student engagement, and adherence to curriculum (just to name a few), the system will provide instructors with actionable feedback. This not only supports professional growth but also ensures that students consistently receive high-quality instruction.

    By bringing these features together into one unified dashboard, the project aims to reduce operational overhead while enabling scalability. Instructors will gain a clear and organized view of their teaching responsibilities, student assignments, and performance metrics. Administrators will benefit from real-time data and analytics that guide decision-making. Most importantly, students will experience smoother, better-matched, and higher-quality instruction—helping CodeSchool continue its mission of making coding a universal literacy.

    In sum, this project is not just a technical upgrade but a step toward reimagining how educational platforms can operate at scale. By blending automation with data-driven insights, it empowers instructors to focus on what they do best—teaching—while the system handles the background logistics. The result is a more efficient, reliable, and impactful platform that supports CodeSchool's global ambitions.
\end{instruction}

\section{Problem definition}
\instruction{
    Currently, managing the lifecycle of classes at CodeSchool is a complex, time-consuming process. Instructors and administrators rely on spreadsheets to schedule classes, assign students, generate Zoom links, and track progress. Evaluations of teaching quality are also carried out manually, with staff watching recorded sessions and giving scores based on observations. While effective at a smaller scale, this system becomes increasingly difficult to manage as CodeSchool grows, leaving room for scheduling conflicts, inefficiencies, and limited insight into teaching quality. These challenges risk slowing down the organization’s ability to expand while maintaining its high standard of student learning outcomes.
}

\section{Social relevance}
\instruction{Describe any societal problem that the project addresses.}

\section{Originality/Novelty}
\instruction{Describe the value of solving the problem. Compare and contrast with any existing solutions.}

\section{CS contribution}
\instruction{Describe the CS component of the project, e.g. the higher level CS courses that contribute to it.}

\section{Scope and Deliverables}
\instruction{Justify the scope of the project with respect to the size of the team and the year long duration. List the foreseeable deliverables.}

\section{Feasibility}
\instruction{List the resources, e.g. datasets, compute resources, software libraries, hardware, required for the project. Mention how you expect to access and utilize them for the project.}

\section{Team dynamics}
\instruction{Justify the suitability of the team members to the project. For example, their relevant courses, projects, internships, or research.}

\section{Tech stack}
\instruction{Write details of the tech stack you will use for this project for e.g. if you are using MERN stack, you can write MongoDB, Express, React and NodeJS etc.}

\section{References}
\instruction{List your references.}

% External advisor undertaking.
\begin{titlepage}

  \newcolumntype{C}[1]{>{\centering\arraybackslash\hspace{0pt}}p{#1}}

  \centerline{\textbf{\ul{Undertaking of Kaavish advisement as an External Supervisor}}}
  \bigskip\bigskip

  I hereby affirm that I have read the project details as described on the preceding pages and agree to undertake advisement of this Kaavish project as an External Supervisor. I understand that this role entails the following.
  \begin{description}
  \item[Meeting] Meeting the project team regularly, at least once every two weeks, for the entire duration of the Kaavish. The meetings may be held remotely if required.
  \item[Advisement] Providing supervision and advice to the team in order to ensure steady progress of the project toward its goals.
  \item[Liaison] Liaising with the Internal Supervisor as required, e.g. to provide feedback or engage in grading.
  \item[Other] Any other task, depending on availability and suitability, relevant to the Kaavish as communicated by the Internal Supervisor or Kaavish Working Group.
  \end{description}

  \bigskip\bigskip\bigskip
  
  \noindent
  \begin{tabular}{@{}lp{.4\textwidth}}
    Name: & \hrulefill\\\\
    Email: & \hrulefill\\\\
    Phone: & \hrulefill\\\\
    Designation: & \hrulefill\\\\
    Affiliation: & \hrulefill\\\\\\
    Signature: & \hrulefill\\
  \end{tabular}
\end{titlepage}


\end{document}

%%% Local Variables:
%%% mode: latex
%%% TeX-master: t
%%% End:
